\documentclass[english]{article}

\usepackage{ae,aecompl}
\usepackage[T1]{fontenc}
\usepackage[latin9]{inputenc}
\usepackage{textcomp}
\usepackage{graphicx}
\usepackage{eurosym}
\usepackage{hyperref}
\usepackage{float}
\usepackage{fancyhdr}
\pagestyle{fancy}
\lhead{\powerenjoy RASD}

\newcommand{\carsharing}{\textit {car sharing }}
\newcommand{\powerenjoy}{\textit{PowerEnjoy }}
\newcommand{\registereduser}{\textit {registered user }}
\newcommand{\registeredusers}{\textit {registered users }}
\newcommand{\staff}{\textit{staff }}
\newcommand{\service}{\textit{service }}
\newcommand{\services }{\textit{services }}
\newcommand{\safearea}{\textit{safe area }}
\newcommand{\safeareas}{\textit{safe areas }}
\newcommand{\powergrid}{\textit{power grid }}
\newcommand{\powergrids}{\textit{power grids }}
\newcommand{\reservation}{\textit{reservation }}
\newcommand{\stand}{\textit{stand }}
\usepackage{tabto}
\usepackage{verbatim}

\makeatletter
\usepackage{babel}


\makeatother

\usepackage{babel}
\begin{document}
\begin{figure}
	\centering
	\includegraphics[scale=0.5]{logo.pdf} 
\end{figure}


\title{\powerenjoy\\
 Requirement Analysis and Specification Document\\
}

\date{A.A 2016/2017}

\author{Erba Alessandro\\
 Leveni Filippo\\
 Lodi Luca}

\maketitle
\pagebreak{}

\tableofcontents{} \pagebreak{}

\section{Introduction}
	\subsection{Purpose }
	
		The purpose of our team is to project \powerenjoy, which is  a platform that will be used to manage a \carsharing \service 
 		 in the city area.\\
 		The platform will allow clients to find and reserve electric cars using mobile app from their self phone, and it will 	also manage the communication of the emergency situations to the proper \staff.\\
  		Furthermore, a particular \powerenjoy feature is the application of certain discounts to its users.\\
 		 \tab \tab This document aims to describe the high-level functionalities that will be offered by \powerenjoy \service. The RASD is intended to be viewed by the stakeholders, that will evaluate the correctness of any assumption and decision in this document.
  
 \subsection{Scope and Problem Description}
 
  The government of a large city would like to introduce an ecological way to travel within the city area. In particular,     uniform distribution of cars in the city is a significant target.\\
  \tab \tab The cars are electric as anticipated and inside them it is available an integrated interface. Through this interface   is possible to communicate with the system (e.g. to call emergency \staff).\\
  \tab \tab Users can choose and reserve a car through a mobile app. The system answers to the request by showing the    time left to the expiration of the \reservation.\\
  The \staff is divided in three categories:\\
   - Field \staff\\
   - Emergency \staff\\
   - Management \staff\\
  \tab \tab Field \staff use the mobile application to inform the system about their availability and to confirm that they are    going to take care of a certain call (e.g. retrieve a broken car, move a vehicle with a low battery or to ensure the uniform    distribution in the city, etc...).\\
  \tab \tab Emergency \staff use the mobile application to answer customers calls and to take care of special situations (e.g. car accident).\\
  \tab \tab Management \staff use the mobile application to handle the \service (e.g. set the \safeareas, fares, discounts,     etc...).\\
  \tab \tab The system guarantees a uniform distribution of cars in the city. The way in which this is guaranteed is explained later in this document.   \\
  \tab \tab A user can take advantage of the "Money saving" option selecting it on the integrated interface in the car. After    that he/she can input his/her final destination and the system provides information about the station where to leave the car to get a discount. This destination is determined to ensure the uniform distribution of cars and depends both on the destination of the user and on the availability of power plugs at the selected station.
	
	\subsection{Goals}
	We have selected a list of objectives that we think \powerenjoy  has to reach:
	\begin{description}
		
		\begin{comment} drivers for \powerenjoy's success \end{comment}
		\item [G.1]{ To allow user to simply access to the \powerenjoy \services. }
		\item [G.2]{ To guarantee an higher probability of finding a car, compared to other \carsharing \services. }
		\item [G.3]{ To guarantee an efficient maintenance and recharge process of the cars, oriented to favorite their fairly distribution over the territory. }
		\begin{comment} user's goals \end{comment}
		\item [G.4]{ To allow users to make a \reservation for a car in advance, and to cancel a \reservation by its expiration time. }
		\item [G.5]{  To transparently charge the user for the \service, minimizing its interaction with payment interfaces. }
		\item [G.6]{  To allow users to drive cars, minimizing their efforts for finding \safeareas where to park.}
		\item [G.7]{ To encourage users' virtuous behaviour  in relation to \service's fairness. }
		\begin{comment} \service efficiency goals \end{comment}
		\item [G.8]{ To create a system that gives the possibility to be monitored and administrated by \powerenjoy's authorized personal. }
		\item [G.9]{  To empower the authorized personal in order to be able to quickly react to emergencies, potentially dangerous situations and users' requests for help.}
		\item [G.10]{ To create a system that provides a programmable interface, in order to reduce the costs of further enhancements/maintenance to the system exploiting internal personal or third party developers.}
	\end{description}

	\subsection{Domain Assumptions }
	We suppose that the following conditions hold in the analysed world:
	\begin{description}
		
		\begin{comment} drivers for \powerenjoy's success \end{comment}
		\item [D.1]{ The geographical area covered by the \service is included in the coverage area of most common mobile communication technologies (3g, 4g) offered by main telecommunications companies. }
		\item [D.2]{ Only correctly \registeredusers can access the functionalities provided by the system. }
		\item [D.3]{ Cars are able to send diagnostic for the manutention process in addition to realtime telemetrics and status of the battery.  }
		\item [D.4]{ Cars can be selectively remotely locked, unlocked, powered on and powered off by authorized personal. }
		\item [D.]{  }
		\item [D.]{  }
		\item [D.]{  }
		\item [D.]{  }
		\item [D.]{  }
		\item [D.]{  }	

	\end{description}

	\subsection{Glossary}
		\subsubsection{Car Sharing}
			\carsharing is a model of car rental where people rent cars for short periods of time, often by the hour.
		\subsubsection{PowerEnjoy}
			\powerenjoy is the \carsharing brand object of this document. 
		\subsubsection{Registered User}
			A \registereduser is a person subscribed to \powerenjoy which can access to the \carsharing through his smart-phone.
		\subsubsection{Staff}
			\staff is the set of people that are \registeredusers but can perform special operations. They are divided in three different categories.
			\begin{itemize}
				\item {Field \staff}
				\begin{itemize}
				\item They own a passepartout for cars that has all kind of issues, for example they can menage cars with low battery.
				\end{itemize}
				\item{Emergency \staff}
				\begin{itemize}
				\item They manage users issues. They are available to communicate with users.
				\end{itemize}
				\item{Management \staff}
				\begin{itemize}
				\item They manage the system, for example they can change system parameters. They can add new \safeareas, modify \carsharing prices.
				\end{itemize}
			\end{itemize}
		\subsubsection{Service}
			The \service is the group of actions that a \registereduser can fulfill trough the mobile.	
		\subsubsection {Safe Area}
			A \safearea is a  geographical place\footnote{defined by a set of GPS positions} on the map in which parking is allowed. \safeareas are saved into the system. A \registereduser can end a \carsharing session only parking used car in a \safearea.
	\subsubsection{Power Grid}
		A \powergrid is the station that charge the battery of an electric car. All \powergrids are displayed in the map on-board \powerenjoy cars. Each \powergrid is in a \safearea. A \registereduser who plugs his car into a \powergrid receives the discount expected within \powerenjoy rules.
	\subsubsection{Reservation}
		A \reservation is the possibility that a \registereduser has to book a car at the latest for one hour.
	\subsubsection{Stand}
		A \stand is the possibility that a \registereduser has to pause his \carsharing. During a \stand the \registereduser can leave the car, than he can re-unlock it through his smart-phone 

\subsection{Further Developments}
	We think that one of the most important things is to release a \service that can be extended and improved. It's trivial that a \service like this must be able to be changed and updated.
	For this purpose we will develop a system open to new extensions. \\
	Here we provide a series of possible extensions that could be necessary in the future.
	E.g.:
	\begin{itemize}
		\item New kind of discount policies can be introduced to the system, for example discounts for new types of users. 
		\item New system policies may be necessary, for example new \carsharing conditions.
		\item New kind of vehicles can be introduced, for example electric scooters and electric bikes.
		\item Support to new payment methods, for example Paypal\textregistered, NFC\footnote{Near Field Communication; NFC devices are used in contactless payment systems, similar to those used in credit cards and electronic ticket smartcards and allow mobile payment to replace/supplement these systems. \href{https://en.wikipedia.org/wiki/Near_field_communication}{NFC on Wikipedia}}.
		\item Develop new \service in order to increase the user experience, for example it may be useful to allow communication among users in order to share a car ride and save money.
	\end{itemize}
	\subsection{Used Tools}
	\begin{itemize}
		\item \LaTeX\\
		\item GitHub\\
	\end{itemize}
\section{Specific Requirements}
	\subsection{Functional Requirements}
	\subsection{Non Functional Requirements}

\section{Scenarios Identifying}
	To allow you to better understand the behavior of this service here we present some possible situations.
	\subsection{Scenario 1}
		Alessandro is planning his weekend. It's his best friend's birthday party and he has to reach a place not covered by public transportation. Alessandro decides to subscribe to the new \powerenjoy \service. He downloads the application on his smartphone, opens the app and clicks on the Sing Up menu. He provides all the information needed to join the \service. In order to do that he sends a scan of his driving license and fills an electronic form containing his personal info (personal details, payment info). Finally he agrees to \service's	"Terms and Conditions'' .
		Now he is waiting for an acknowledgment of registration. After a while he receives an email reporting his successfully subscription to the \service. Alessandro can now Log In in the application as a \registereduser.
	\subsection{Scenario 2}
		Ariana feels smart, last week she has rent a \powerenjoy but she has run out of money in her prepaid card, so the system has locked her account. She owns 2 different prepaid credit cards and she tries to re-subscribe to the \service without settle her debt. Ariana provides all the information for subscribing to the \service and sends it. Her application request is denied.
	\subsection{Scenario 3}
		It's Saturday night and Alessandro is a newbie of \powerenjoy \service, he tries to reserve his car, he log in in the app and turns on the GPS of his smartphone. In a few seconds a map appears on the screen. Here are displayed all the available cars around him. Clicking on a car he can see it's autonomy. He choses the nearest car to his house and reserves it clicking on the ``Reserve Now'' button. From this moment he has an hour before it's \reservation expires. He gets ready for the night and than walks to the car. Once arrived he opens the 	\powerenjoy app and goes to the ``Here I am'' menu. The system checks Alessandro's GPS Position and if he is close enough to the car unlocks it. Alessandro is very happy and gets in the car.
	\subsection{Scenario 4}
		It's Monday morning and Alice is getting ready for her work day. Alice is a \registereduser of \powerenjoy . She is late and won't take the subway today. She reserves a car near her apartment. While she is getting ready for the day, she receives a call from Melissa. Melissa is an Alice working colleague. Melissa is gonna picking Alice up. Alice cancels hers \reservation and pays a fee of 0.25\euro because it is passed less than a quarter of an hour from his \reservation time.
	\subsection{Scenario 5}
		Marcus has to carry his son Mario to school. He reserves a \powerenjoy. While they are going downstairs Mario falls down. He loses consciousness and Marcus calls an Ambulance. In this situation Marcus completely forgets to cancel his \reservation, so after an hour 1\euro is charged	on his credit card.
	\subsection{Scenario 6}
		Antony is moving in the subway from Duomo to Sesto Primo Maggio Fs. It's Friday 21th of October and in Sesto Primo Maggio Fs there are not trains to reach Lissone because there is an usual Trenord strike. Fortunately he is a \powerenjoy user. Antony Logs In in the app and selects the manual input car research. Typing on his smartphone ``Sesto San Giovanni'' he easily checks all the available  cars. He choses a car and reserves it. He is very happy because the "Money Saving Option'' guarantees him an high probable possibility to find a car in each \powerenjoy \safearea. Once in Sesto San Giovanni he unlocks his car while all the other commuter can only hope to reach their home within the day after.
	\subsection{Scenario 7}
		It's raining. Filippo is going towards the \powerenjoy car just reserved when he meets Alessandro and Luca, two his classmates. They have to go home, Filippo offers to give a lift to both for the station and they accept. Once into the car, the system notes that there is two passengers. Arrived at the station Filippo stops the car and his friends leave the veichle. At this time the system store on its memory that on that stretch of road traveled it will have to apply the discount of 10\%.
	\subsection{Scenario 8}
		Andrew has to carry his mom to the airport but remembers that his car is being repaired. He thinks that the best choice is to reserve a car with the \powerenjoy app. After that they go towards the vehicle and, once reached it, Andrew helps his mom to load the heavy suitcase in the back seat. Now everything is ready to go. On the way by the airport on the screen of the car appears a message that explain to the driver that he finished the money on his credit card, so he can finish this trip but his account will be blocked until he will pay the debts towards the company. Arrived at the airport the mom goes out from the car with her suitcase and says goodbye to her son. Once at home Andrew has to pay the full rent because the system notice that the passenger is only one.
	\subsection{Scenario 9}
		Paul is late for his first day of teaching. Go by foot will not allow him to arrive on time. He decides to reserve a car with his \powerenjoy app. He opens it and checks the available cars. Unfortunately he notices that the only nearby car have 15\% of battery full and the workplace is 4KM from the nearest \powergrid station, outside a \safearea. The workplace is just 10 minutes by car from its position and the time is running out, so 	he reserves it. Left the \safearea, the car advice Paul by the screen that he should 	park in a \safearea, otherwise will be charged recovery costs. He's too much late and ignores the advice. The system will charge 30\% more on this ride because the car is left with less 20\% of the battery's charge and at more that 3KM from the nearest \powergrid station, furthermore Paul will has to pay recovery costs.
	\subsection{Scenario 10}
		Rodolfo has to meet her parents in their downtown manor. He could go by feet but it's a very cold winter day and prefers to reserve a warmer \powerenjoy car. Looking at the cars around him, luckily notices that the nearest car have the battery completely charge! He is clever and kind, so decide not to go directly to his parent's house but to reach the closest station to it in order to plug the car into the \powergrid. Rodolfo reaches the car, get in and starts driving. Arrived at the station, the car's battery level is 90\%. For this journey he will receive a total of 50\% discount: 20\% from the level of the battery ( more than 50\% ) and 30\% from his kindness ( car plugged into a \powergrid ). A very lucky day for Rodolfo!
	\subsection{Scenario 11}
		Amelia is a good mom and be careful to save money for her family, in fact she went to the supermarket by foot. After buying the necessary things for the house, she notices that the bags are too heavy to walk back. She remembers that she is subscribed to the \powerenjoy \service. She logs in the app, reserves a car. Than she picks up the car, selects the destination and chooses the ``money saving option''. Although this unforeseen she can save some money. The system computes the destination in order to get the discount. Once arrived to the destination she locks the car and plugs it into the \powergrid. The system applies a discount of 30\% on her ride.
	\subsection{Scenario 12}
		John is going to the gym using the \powerenjoy \service, when he notices that the car's battery level is critically low, but he decides to go on without takes care of it. At certain point the car informs John that he should go near the roadside because the car will stops shortly. John is a bully and decides to ignore the warning. Taking care of the user's safety the car stops as expected, but John is in the middle of the road! Never having been in a similar situation he decides to call the assistance using the embedded screen on the car. An operator takes the call and explains John that a field \staff is coming on site and until the car is not parked safely he will continue to pay. A very expensive day for our John.		
	
\section{UML Models And Use Cases}
	\subsection{Use Cases Diagram}
	\subsection{Actors Identifying}
	\subsection{Use Cases}

	\subsection{Class Diagram}
\section{External Interfaces}
\section{Alloy Model}
\section{Hours of Work}

\end{document}